\section{Notation, Terminology, and Background}
\label{sec:background}

We fix some standard notation across the paper.

Coinduction based on greatest fix points of monotone functions. We denote
greatest fix points as \gfp{}, and an arbitrary monotone function as \bfun{}.

Paco is parameterized coinduction; tower is the tower induction approach of
Schäfer and Smolka~\cite{schafer-smolka} adapted by Pous~\cite{pous-companion}.

\subsection{Bisimilarity}

Up-to techniques are functions that satisfy this law:
\gap{State the compatibility law. $f \circ b \leq b \circ f$.}

This lets us write an easier coinductive proof.
\gap{The easier proof. Smallest example that shows it.}

\subsection{Interaction Trees}

An Interaction Tree (ITrees) is a coinductive data structure for representing
impure and diverging computation inside a proof assistant (which is necessarily
pure and terminating). The full details of ITrees are too complex to cover in
full in this paper; the curious reader may inspect the original Interaction
Trees paper~\cite{itrees}.

It suffices to say that the definitions of strong and weak bisimilarity over
itrees, and the proofs depending on those definitions, rely heavily on
coinductive machinery. Interaction Trees was originally constructed using paco,
as it enabled \gap{\ldots}

\subsection{Problems with paco}

But paco has problems and/or warts:
\begin{itemize}
  \item Many tactics, and many duplicates across arities.
  \item Hard to read and write compatibility proofs: \gap{example}
\end{itemize}

Theory: enhanced coinduction, with its fewer tactics, all-arity story, and
uniform compatibility proofs, will be a simpler theory to underly itrees and
will lead to shorter proofs.

This is important because abstractions leak, as mentioned above. When they do,
we want an easy-to-use theory!
